\documentclass[12pt]{article}
\usepackage[T1]{fontenc}
\usepackage[utf8]{inputenc}
\usepackage{lmodern}
\usepackage{textcomp}
\usepackage{enumitem}
\usepackage{geometry}
\usepackage{graphicx}
\usepackage{amsmath, amssymb}
\geometry{margin=1in}
\begin{document}
\begin{center}
Physics - Undergraduate Level Assessment 24 \
Total Marks: 40 \
Answer all questions.
\end{center}

\section*{Multiple Choice (10 marks)}
\begin{enumerate}[label=\arabic*.]
\item In the diamond structure of elemental carbon, the nearest neighbors of each C atom lie at the corners of a
\begin{enumerate}[label=(\alph*)]
    \item square
    \item hexagon
    \item cube
    \item tetrahedron
\end{enumerate}
\item Which of the following statements about bosons and/or fermions is true?
\begin{enumerate}[label=(\alph*)]
    \item Bosons have symmetric wave functions and obey the Pauli exclusion principle.
    \item Bosons have antisymmetric wave functions and do not obey the Pauli exclusion principle.
    \item Fermions have symmetric wave functions and obey the Pauli exclusion principle.
    \item Fermions have antisymmetric wave functions and obey the Pauli exclusion principle.
\end{enumerate}
\item Which of the following is true about any system that undergoes a reversible thermodynamic process?
\begin{enumerate}[label=(\alph*)]
    \item There are no changes in the internal energy of the system.
    \item The temperature of the system remains constant during the process.
    \item The entropy of the system and its environment remains unchanged.
    \item The entropy of the system and its environment must increase.
\end{enumerate}
\item By definition, the electric displacement current through a surface S is proportional to the
\begin{enumerate}[label=(\alph*)]
    \item rate of change of the electric flux through S
    \item electric flux through S
    \item time integral of the magnetic flux through S
    \item rate of change of the magnetic flux through S
\end{enumerate}
\item By definition, the electric displacement current through a surface S is proportional to the
\begin{enumerate}[label=(\alph*)]
    \item magnetic flux through S
    \item rate of change of the magnetic flux through S
    \item time integral of the magnetic flux through S
    \item rate of change of the electric flux through S
\end{enumerate}
\end{enumerate}

\section*{True / False (10 marks)}
\textit{Answer True or False}
\begin{enumerate}[label=\arabic*.]
\setcounter{enumi}{5}
\item True or False: In the observer's reference frame, it takes 2.5 ns for a meter stick moving at 0.8 c to pass the observer.
\item True or False: The angular magnification of the telescope is 5, which is calculated using the separation distance of the lenses.
\item True or False: An electron in the l = 2 state can have only three allowed values of the magnetic quantum number m\_l.
\item True or False: The sign of the charge carriers in a doped semiconductor can be determined by measuring the Hall coefficient.
\item True or False: According to BCS theory, the attraction between Cooper pairs in a superconductor is due to the weak nuclear force.
\end{enumerate}

\section*{Long Form Response (20 marks)}
\textit{Answer each question with a clear and complete explanation.}
\begin{enumerate}[label=\arabic*.]
\setcounter{enumi}{10}
\item Discuss the conditions under which a beam of electrons can diffract when interacting with a crystal surface, and analyze the approximate kinetic energy required to observe the resulting scattering pattern.
\item Discuss the process and calculations involved in determining the final voltage across a charged capacitor when two initially uncharged capacitors are connected in series to it.
\end{enumerate}

\vfill
\noindent\textit{End of Paper}
\end{document}